\chapter{Windows I: Cross-Compiling from Linux}
\label{ch:windows-cross-compiling}
\label{ch:windows-wine}
\index{cross-compilation}\index{MinGW-w64}\index{Windows}

CNA's evidenced Windows engineering loop starts on Linux. CMake selects a MinGW target,
builds target-triple-specific SDL dependencies, links a PE executable, stages its runtime
DLLs, and teaches CTest to invoke that executable through Wine. Each step has an independent
failure mode; a successful compiler probe is not yet a runnable Windows program.

\section{Two MinGW toolchains}

\texttt{cmake/toolchains/mingw-w64.cmake} targets Windows x86\_64 when the matching compiler
exists and otherwise falls back to the i686 triple. It sets C, C\texttt{++}, and resource
compilers and constrains CMake package, library, and header searches to the cross root while
allowing build-host programs to remain native.

\begin{lstlisting}[style=cnashell]
cmake -B build-windows \
  -DCMAKE_TOOLCHAIN_FILE=cmake/toolchains/mingw-w64.cmake \
  -DCNA_GRAPHICS_RENDERER=SDL_RENDERER
cmake --build build-windows
\end{lstlisting}

The fallback is not a fully self-consistent 32-bit identity: the file leaves
\texttt{CMAKE\_SYSTEM\_PROCESSOR} at \texttt{x86\_64} even after selecting the i686 compiler
and search root. It can rescue compiler discovery, but it must not be treated as equivalent to
the dedicated i686 lane when processor-keyed caches or ABI gates matter.

The separate \texttt{mingw-w64-i686.cmake} is not merely a convenience alias. GLIDE's native
ABI represents a window handle as a 32-bit \texttt{FxU32}, and the available Glide DLL is x86.
That toolchain fixes the processor and all compiler names to i686. As the final section shows,
it can validate the renderer boundary but cannot link the whole framework.

There is no CNA-owned Android, Emscripten, or Apple toolchain file. Those paths use the NDK or
emsdk toolchain externally; Apple requires its native environment. The two checked-in files are
specifically the Windows-from-Linux contract.

\section{Vendored SDL is built for the target, not reused from the host}

Before renderer selection, CNA unconditionally calls \texttt{cna\_configure\_vendored\_sdl()}.
Its helper performs a nested configure, build, and install of SDL3, SDL3\_image, and
SDL3\_mixer, forwarding the parent's toolchain. A MinGW build therefore receives Windows
libraries and DLLs rather than accidentally linking the Linux host's SDL.

The cache directory is keyed by system and processor:

\begin{lstlisting}
.sdl-prebuilt-Linux-x86_64/
.sdl-prebuilt-Windows-x86_64/
.sdl-prebuilt-Android-aarch64/   # expected by one demo, absent at audit
.sdl-prebuilt-emscripten/       # absent at audit
\end{lstlisting}

The audited machine contains the first two only. The Windows cache records the MinGW toolchain
and GNU import libraries such as \texttt{libSDL3.dll.a}; it has no MSVC \texttt{.lib} or
\texttt{.pdb}. This is direct evidence of cross-compilation, not native Visual Studio output.

Configure-time dependency builds make first configuration heavier and more stateful, but they
also keep the target triple explicit. A stale cache from another ABI should not be renamed into
place; it must be regenerated under the intended toolchain.

\section{Runtime staging is part of a successful build}

A PE executable that links against an import library still needs the corresponding DLL at
runtime. CNA's post-build helpers copy SDL3, SDL3\_image, and SDL3\_mixer beside Windows
targets. MinGW test targets link \texttt{libgcc} and \texttt{libstdc++} statically, then stage
\texttt{libwinpthread-1.dll}. Application targets with large RTTI graphs may use the dynamic
C\texttt{++} runtime; a separate helper locates and copies
\texttt{libgcc\_s\_seh-1.dll}, \texttt{libstdc++-6.dll}, and the threading DLL.

Each helper asks the selected compiler where its runtime file lives and falls back to the
compiler's bin directory where appropriate. Failure is a warning that the executable may not
run on a clean prefix, rather than a link failure on the build host. Verification must therefore
inspect or execute the staged directory, not stop at ``target built.''

Renderer-specific dependencies add another layer. DIRECTX8 needs a D8VK
\texttt{d3d8.dll.a}; Direct3D renderers need the corresponding Windows SDK import libraries;
GLIDE needs an external \texttt{glide3x.dll} only at execution. These are not supplied by the
generic MinGW runtime helper.

\section{CTest must know that PE files are not host executables}

GoogleTest discovery in \texttt{PRE\_TEST} mode executes \texttt{CnaTests.exe} to enumerate
cases. Without an emulator property, Linux attempts to run the PE file natively and discovery
fails before any test. CNA sets CMake's \texttt{CROSSCOMPILING\_EMULATOR} property on
\texttt{CnaTests} for DirectX1/2/3/5/6/7/8/9/10/11/12 and Direct2D, selecting the matching Wine
wrapper.

Discovery and device-free unit tests set each wrapper's documented skip-gate variable. A bare
\texttt{--gtest\_list\_tests} call creates no Direct3D device and cannot emit a DXVK marker; an
engagement gate would falsely reject it. Renderer smoke examples instead put the wrapper in the
registered command because those executables really do create the device and should prove the
runtime engaged.

\begin{lstlisting}
# Conceptual CMake behavior
set_target_properties(CnaTests PROPERTIES
  CROSSCOMPILING_EMULATOR
    "env;CNA_D3D9_SKIP_DXVK_GATE=1;scripts/run-wine-dxvk9.sh")
\end{lstlisting}

The strict public-API compile check has its own emulator wiring. These two sites are the bridge
that makes ordinary \texttt{ctest} understand the cross-built artifacts; per-example commands
remain explicit.

\section{Native MSVC is a separate, unproven lane}

Two Windows workflows target \texttt{windows-latest}, activate MSVC, use Ninja, and include
tests. Both are manual \texttt{workflow\_dispatch} workflows. The repository cites no Actions
run ID for them, and CNA production code has essentially no hand-written MSVC conditional
logic: its only \texttt{\_MSC\_VER} sites are generated Sokol shader output.

That is not proof that MSVC fails. It means the evidence in this repository establishes
MinGW-plus-Wine, not a native Microsoft toolchain or genuine Windows execution. Workflow YAML
is executable intent; a recorded run is execution evidence.

\section{The i686 wall}

The 32-bit GLIDE lane reaches farther than a syntax sketch. The renderer passes
\texttt{-fsyntax-only}, and a fake DLL client verifies all 39 expected exports under Wine. The
full CNA executable cannot link because sharp-runtime's \cnaclass{Int128} and
\cnaclass{Decimal} require GCC/Clang \texttt{\_\_int128}, which
\texttt{i686-w64-mingw32-g++} does not provide. Sharp-runtime explicitly declines to hand-roll
128-bit arithmetic solely for MSVC/32-bit coverage.

GLIDE therefore has source, configuration, partial cross-compile, and an ABI probe, but no full
artifact. The renderer's own syntax can be correct while an unconditional foundation dependency
makes the product architecture unreachable. This is exactly why platform status belongs to a
layered evidence vector rather than a single build-support checkbox.
